\documentclass{report}
\usepackage[english]{babel}

\usepackage[most,many,breakable]{tcolorbox}
\usepackage{xcolor}

\definecolor{defBoxBorder}{HTML}{395144}
\newtcolorbox{defBox}{colback=white,colframe=defBoxBorder,arc=3pt, boxrule=0.5pt, drop fuzzy shadow, title=Definition}
\definecolor{thBoxBorder}{HTML}{AC8441}
\newtcolorbox{thBox}{colback=white,colframe=thBoxBorder,arc=3pt, boxrule=0.5pt, drop fuzzy shadow, title=Theorem}
\definecolor{noteBoxBorder}{HTML}{4E6C50}
\newtcolorbox{noteBox}{colback=white,colframe=noteBoxBorder,arc=3pt, boxrule=0.5pt, drop fuzzy shadow, title=Note}
\definecolor{axBoxBorder}{HTML}{AA5656}
\newtcolorbox{axBox}{colback=white,colframe=axBoxBorder,arc=3pt, boxrule=0.5pt, drop fuzzy shadow, title=Axiom/Postulate}
\definecolor{corBoxBorder}{HTML}{8B7E74}
\newtcolorbox{corBox}{colback=white,colframe=corBoxBorder,arc=3pt, boxrule=0.5pt, drop fuzzy shadow, title=Corollary}
\definecolor{lemBoxBorder}{HTML}{B99B6B}
\newtcolorbox{lemBox}{colback=white,colframe=lemBoxBorder,arc=3pt, boxrule=0.5pt, drop fuzzy shadow, title=Lemma}


\usepackage[spanish]{babel}
\usepackage[utf8x]{inputenc}
\usepackage{amsmath}
\usepackage{graphicx}
\usepackage[colorinlistoftodos]{todonotes}
\usepackage{enumitem}
\usepackage{listings}
\usepackage{verbatim}
\usepackage{eurosym}
\usepackage[export]{adjustbox}
\usepackage{amssymb}
\usepackage{bussproofs}
\usepackage{amsmath}
\usepackage{tikz}
\usepackage{xcolor}
\usepackage{listings}
\usepackage{titletoc}
\usepackage{hyperref}

\hypersetup{
  colorlinks=true,
  linkcolor=black,
  urlcolor=blue,
  citecolor=black
}

\newcommand{\coverPage}[6]{%
%----------------------------------------------------------------------------------------
%	COVER START
%----------------------------------------------------------------------------------------
\begin{titlepage}

    \newcommand{\HRule}{\rule{\linewidth}{0.5mm}}
    \newcommand{\department}{#1}
    \newcommand{\course}{#2}
    \newcommand{\titleValue}{#3}
    \newcommand{\subtitleValue}{#4}
    \newcommand{\authorName}{#5}

    \center

    %----------------------------------------------------------------------------------------
    %	HEADER
    %----------------------------------------------------------------------------------------
    \includegraphics{images/logo_usa.png}
    \vspace{0.5cm}
    \textsc{\Large \department}\\[0.5cm]
    \textsc{\Large \course}\\[0.5cm]
    \vfill

    %----------------------------------------------------------------------------------------
    %	TITLE
    %----------------------------------------------------------------------------------------

    \HRule\\
    \Huge
    \textbf{\titleValue}\\[0.5cm]
    \Large
    \textbf{\subtitleValue}\\
    \HRule\\[0.5cm]

    %----------------------------------------------------------------------------------------
    %	AUTHOR AND DATE
    %----------------------------------------------------------------------------------------

    \vfill
    \Large
    \textit{\authorName}\\
    {\large \today}\\[2cm]

\end{titlepage}
%----------------------------------------------------------------------------------------
%	COVER END
%----------------------------------------------------------------------------------------
}

\begin{document}
    \coverPage{ Mathematics }{ Integral Calculus and Series }{ Integrability }{  }{ Alexander Mendoza }{\today}
    \tableofcontents

    \pagebreak
    \chapter{ Integrability }
    % Add a description of the chapter on weekend revision

    % Add an introduction to the relation between integral and derivatives
    To define the integral we will first need some definitions.
    \subsection*{Inferior and Superior Sums}
    We will use both the superior and inferior sums to define the concept of an integral.

    \begin{defBox}
        \textit{\textbf{Partition of an integral.}} Let $a, b$ be real numbers such that $a<b$ and let $P$ be the following set

        $$P = \{x_0, x_1, \dots , x_n\}$$ with $a=x_0<x_1<\cdots<x_n=b$, then we say that $P$ is a partition of $[a,b]$
    \end{defBox}

    % Insert graphic
    Figure 1, represents graphically the concept of partition.

    \begin{defBox}
        Let $f$ be a function bounded on $[a,b]$ and let $P = \{x_0, x_1, \dots x_n\}$ be a partition of $[a,b]$. then

        The minimum of a function in $[a,b]$ is defined as
        $$m_i := inf\{f(x) | x_{i-1} \leq x \leq x_i\}$$
        And the maximum of a function in $[a,b]$ is defined as
        $$M_i := sup\{f(x) | x_{i-1} \leq x \leq x_i\}$$
    \end{defBox}

    % Add examples

    With this we can now define the upper and lower sums.

    \begin{defBox}
        Suppose $f$ is bounded on $[a,b]$, let $P = \{t_0,\dots,t_n\}$ be a partition of $[a,b]$ and let

        $$m_i = inf\{f(x) | x_{i-1} \leq x \leq x_i\}$$
        $$M_i = sup\{f(x) | x_{i-1} \leq x \leq x_i\}$$

        Then we define the lower sum of $f$ for $P$, denoted by $L(f, P)$, as

        $$L(f, P) = \sum_{i=1}^{n}m_i(t_{i} - t_{i-1})$$

        and the upper sum of $f$ for $P$, denoted by $U(f, P)$, is defined as

        $$U(f, P) = \sum_{i=1}^{n}M_i(t_{i} - t_{i-1})$$
    \end{defBox}

    % add examples

    Let's now see which relation the lower and upper sum have. Our intuition may tell us that given $f$ bounded over $[a,b]$ and $P$ a partition of $[a,b]$, $L(f, P) \leq U(f, P)$, in fact this is true for any pair of distinct partitions. But, let's review first a lemma to demonstrate this fact.

    \begin{lemBox}
        Let $f$ be a bounded function and $P, Q$ be partitions of $[a,b]$. If $Q \subseteq P$, then

        $$L(f, Q) \leq L(f, P) $$
        $$U(f, P) \leq U(f, Q) $$
    \end{lemBox}

    \begin{ideaBox}
        The idea for this proof is to first demonstrate that the lemma is true for a partition that contains only one more element, and then construct a succession of partitions, where each new partition will have one more element than the partition before it, and proof the general case.
        % Expand more on this idea
    \end{ideaBox}

    \textit{\textbf{Proof.}} Let $f$ be a bounded function and $P, Q$ be partitions of $[a,b]$ where $Q \subseteq P$ and $P$ contains only one more element than $Q$, namely $u$. We can represent $P$ and $Q$ as

    $$P = \{x_0,x_1,\dots,x_{k-1},x_k,\dots,x_n\}$$
    $$Q = \{x_0,x_1,\dots,x_{k-1},u,x_k,\dots,x_n\}$$

    Then by definition we have

    \begin{align*}
        L(f,P) &= \sum_{i=1}^{n} m_i(x_i-x_{i-1})\\
        &= \sum_{i=1}^{k-1}m_i(x_i-x_{i-1}) + m_k(x_k-x_{k-1}) + \sum_{i=k+1}^{n}m_i(x_i-x_{i-1})
    \end{align*}

    Similarly

    \begin{align*}
        L(f,Q) &= \sum_{i=1}^{n} m_i(x_i-x_{i-1})\\
        &= \sum_{i=1}^{k-1}m_i(x_i-x_{i-1}) + m'(u-x_{k-1}) + m''(x_k-u) + \sum_{i=k+1}^{n}m_i(x_i-x_{i-1})
    \end{align*}

    % COMPLETE THE PROOF

    With this we can state a theorem that will allow us to then define the integral of a function. Before providing the theorem, a useful observation is that, given any partition $P$
    $$L(f, P) \leq U(f, P)$$
    We know this because
    $$L(f, P) = \sum_{i=1}^{n}m_i(t_{i} - t_{i-1})$$
    $$U(f, P) = \sum_{i=1}^{n}M_i(t_{i} - t_{i-1})$$
    and by definition, for any $i$:
    \begin{align*}
        m_i &\leq M_i\\
        m_i(t_{i} - t_{i-1}) &\leq M_i(t_{i} - t_{i-1})
    \end{align*}

    \begin{thBox}
        Let $P_1, P_2$ be partitions of an interval $[a,b]$ and let $f$ be bounded over $[a,b]$, then it's true that

        $$L(f, P_1) \leq U(f, P_2)$$
    \end{thBox}

    \textit{\textbf{Proof.}} Let $P = P_1 \cup P_2$, then by definition of union, $P_1 \subseteq P$ and $P_2 \subseteq P$, and by the previous lemma:

    $$L(f, P_1) \leq L(f, P)$$
    $$U(f, P) \leq U(f, P_2)$$
    By the afore mentioned observation, we have
    $$L(f, P) \leq U(f, P)$$
    And all together
    $$L(f, P_1) \leq L(f, P) \leq U(f, P) \leq U(f, P_2)$$
    Finally, by transitivity, we therefore have
    $$L(f, P_1) \leq U(f, P_2)$$
    Concluding our proof.

    % EXPAND ON WHY THIS IS IMPORTANT

    With this, we can now define the integral of a function.

    \begin{defBox}
        Let $f$ be a function bounded on $[a,b]$, $f$ is said to be integrable on $[a,b]$ if


        \begin{align*}
            sup\{L(f,P) | P \text{ is a partition of } [a,b]\}\\
            = inf\{U(f,P) | P \text{ is a partition of } [a,b]\}
        \end{align*}

        That common number is what we call the integral of $f$ over $[a,b]$ and is denoted by

        $$\int_{b}^{a}f$$

        And is read "the integral of $f$ over $[a,b]$"
    \end{defBox}

    % INSERT EXAMPLES
\end{document}
